% vision_pipeline
\section{视觉感知流水线}

网球机器人的视觉感知从一帧图像开始，到网球在画面中的坐标结束。相机取来一帧，解码、缩放与色彩转换把它整理成检测网络要求的输入，NPU 推理给出一组候选框，后处理在这些框里筛出网球并算出它的位置。这条流水线的每一个阶段都有两条落点，既可以留在 CPU 上由通用核逐像素完成，也可以交给 RK3588 上对应的专用硬件单元。JPEG 解码、二维缩放与色彩转换、张量推理这三步在 CPU 上各自的开销都与一次推理相当，若这三步均由通用核承担，感知与控制将争抢同一组核，流水线随之受阻。将这三步逐级交由 JPU、RGA、NPU 三个专用单元处理后，通用核方可释放算力用于决策与运动控制。

\figphl{../figures/out/Fvision.png}{视觉感知流水线：相机取流 → JPU 解码 → RGA 缩放与色彩转换 → NPU 推理 → 后处理}{fig:vision-arch}

图~\ref{fig:vision-arch} 给出这条流水线的分级结构。其中相机的 USB/xHCI 主机栈与 NPU 的 \texttt{/dev/card1} 节点为既有能力，JPU 驱动、RGA 驱动、dma-heap 分配器与 NPU 的 GEM 缓冲修复为本队新增。以下逐级说明各硬件单元与对应的工作。

\subsection{流水线各阶段}

\paragraph{相机}
流水线的第一级是接在 RK3588 USB3 端口上的相机。它符合 USB 视频类（UVC）规范，经 Starry 既有的 xHCI 主机栈以等时传输取流，无需专用驱动。这只相机交付两种格式，YUYV 是未压缩的 4:2:2 原始像素，取来只需一次色彩转换；MJPEG 把视频的每一帧各自压成一张独立的 JPEG，取来还要先解码。两种格式后续都进入 RGA 与 NPU，区别只在前端是否多一步解码。这一级复用既有主机栈，取流通路本身不引入新驱动。xHCI 端点上下文的 \texttt{Interval} 字段决定控制器多久为等时端点取一次数据，它并不直接等于描述符里的 \texttt{bInterval}，须按端点类型与连接速度换算。对高速及更高速度的等时端点，正确换算是 \texttt{Interval} 取 \texttt{bInterval} 减一，含义为每 $2^{\texttt{Interval}}$ 个 125\,$\mu$s 微帧服务一次；\ours{\texttt{usb-host} 的 \texttt{calculate\_xhci\_interval} 按速度与端点类型分类固定这一换算}。若为简化实现而以一个固定下界将该值截断至 1，\texttt{bInterval} 为 1 的高带宽端点会被编码成每两个微帧才服务一次，可用带宽当场减半，相机内部缓冲随即溢出，等时帧被成片截断。同一段端点配置还显式把等时端点的错误计数置零（等时传输本不重传），并按每微帧包数算出突发大小与一次服务区间的最大载荷（max ESIT payload）一并写入上下文，控制器才会按相机申报的带宽真正预约总线。对实到字节短于申报长度的等时帧，驱动逐包归类为足额、偏短或为零，并与缺失完成事件的包数分开统计，从而把事件投递丢失与控制器侧限速或 FIFO 截断两类成因区分开，这是在板上诊断一条被截断的等时输入流时唯一能凭代码坐实的判据。

\paragraph{JPU}
MJPEG 帧的解码交给 JPU，即 RK3588 集成的硬件 JPEG 解码单元（设备树节点 \texttt{jpegd@fdb90000}）。在 CPU 上软解一帧约需 25\,ms，几乎与一次推理相当；JPU 在硬件中完成熵解码与反变换，直接将一张 JPEG 还原为 \texttt{NV12} 帧并写入物理内存，从而将这一步移出计算路径。Starry 此前没有这块硬件的驱动，现已\ours{新增其平台无关的驱动核心与 \texttt{/dev/mpp\_service} 设备节点}，对齐瑞芯微 MPP 库的会话协议，使未经改动的 MPP 程序照常驱动这块解码器。MPP 与内核之间是一串按字节对齐、经 ioctl 传入的请求记录，每条为一个 24 字节结构，带命令号、标志位与一个指向各自负载的用户态指针，多条以 \texttt{MULTI\_MSG} 标志串成一批至 \texttt{LAST\_MSG} 为止；会话先以 \texttt{PROBE\_HW\_SUPPORT} 询问本节点支持的编解码并据实回报只支持 JPEG 解码，再以 \texttt{QUERY\_HW\_ID} 读硬件版本、以 \texttt{INIT\_CLIENT\_TYPE} 把会话绑定到 JPEG 解码客户端类型，随后 \texttt{SET\_REG\_WRITE} 送入整组寄存器、\texttt{SET\_REG\_READ} 声明完成后回读的寄存器窗口，最后 \texttt{POLL\_HW\_FINISH} 触发一次提交并阻塞到这一帧解码完成。提交前节点按厂商 \texttt{trans\_tbl\_jpgdec} 给定的位置逐一扫描地址槽，把 dma-buf 文件描述符经共享的 \texttt{/dev/dma\_heap} 解析为其连续缓冲的物理基址、再加声明的偏移写回寄存器，量化表、霍夫曼最小码表与霍夫曼值表共处一段表缓冲、分别落在偏移 384 与 704 字节处。设备初始化时把前置的 IOMMU 置为直通：其寄存器在同一寄存器页偏移 \texttt{0x480} 处，先写入强制复位命令清掉前一个系统可能遗留的页表，再把页表基址寄存器 \texttt{RK\_MMU\_DTE\_ADDR} 清零、写入停用分页命令 \texttt{RK\_MMU\_CMD\_DISABLE\_PAGING} 并屏蔽全部中断，此后引擎按物理地址直接读写内存。\ours{运行时把整组寄存器写入硬件时跳过只读的 ID 寄存器，并把带启动位的控制寄存器 \texttt{SWREG1} 留到最后一个写}，从而保证几何、格式、表长与五个地址槽都就位之后引擎才被拉起，绝不会在输入尚未配齐时就开始取数。开发板上单帧解码中位约 1\,ms，输出与 Linux 逐字节一致。该驱动以合并请求 \href{https://github.com/rcore-os/tgoskits/pull/1456}{\#1456} 提交至上游，已合并。

\paragraph{RGA}
缩放与色彩转换交给 RGA，即 RK3588 集成的二维栅格图形加速器，它在硬件中完成色彩空间转换、缩放与图层拼接。相机交付的是 YUV 排布的像素，而检测网络要求连续的 RGB 平面，尺寸还需缩放到模型的输入分辨率；在 CPU 上逐像素做这件事，一帧的开销与一次推理相当。Starry 此前没有这条通路，现已\ours{新增平台无关的 RGA2 驱动核心与 \texttt{/dev/rga} 设备节点}，按 librga 的真实二进制接口实现 \texttt{rga\_req} 结构体布局与 \texttt{RGA\_BLIT\_SYNC} 约定，使既有的 librga 二进制无需改动即可驱动 RGA2 核。\texttt{rga\_req} 在 LP64 下恰为 504 字节，内嵌三份各 56 字节的图像描述符，缓冲导入用的 \texttt{rga\_external\_buffer} 为 288 字节，这些尺寸以编译期断言固定，任何字段增删都会在编译时暴露；早期镜像漏掉 \texttt{rotate\_mode}、\texttt{rd\_mode} 等较新字段，使其后所有字段整体偏移，引擎据此取到错位的旋转矩阵与地址。librga 在打开设备时查询驱动版本与硬件版本以选择结构体布局，因此设备节点须如实报告 MultiRGA v1.3.1 与 RK3588 上 RGA2 核的版本号，librga 才会按这一布局组包并把请求投向承载 RGA2 的那个核。RGA2 的执行时序对寄存器写入顺序敏感：引擎须先被置入主模式，命令起始信号才会触发一次命令 DMA 取指，且表示全部命令完成的状态位只有在对应中断使能于起始之前被置位时才会锁存，最初省略这一步会使引擎执行完毕后完成位仍为零，导致轮询持续超时；\ours{据厂商时序在主模式与起始之间补上对中断使能位的读改写}后完成位正常锁存，既有的忙等轮询即可观察到它，无需真正布线中断。缩放系数以 16.16 定点表示，缩小时正确系数是目标维度比源维度、其值不超过一并写入低半字，最初误用其倒数使源指针每步跨过有效区域，扫描器随即停在忙状态约 50\,ms 不再前进，按缩小与放大两种情形分别依厂商参数计算后缩放回到一次提交即完成。该驱动以合并请求 \href{https://github.com/rcore-os/tgoskits/pull/1388}{\#1388} 提交至上游，已合并。

\paragraph{NPU}
整理好的一帧送进检测网络，承担这一步的是 RK3588 片上集成的神经网络加速器 NPU。运行时以用户态共享库 \texttt{librknnrt.so} 加载 RKNN 模型、提交一帧、取回结果，经 \texttt{/dev/card1} 这个 DRM 字符设备下达到内核。这个节点及其 RKNPU 驱动本已存在，但要承载一帧接一帧的连续推理，GEM 缓冲对象在映射与销毁两处的行为须被修正。映射时若不正确处理这段区域的缺页，第一次访问就触发无法收场的异常，紧随其后的推理取不到输入；销毁时若引用计数处理有误，被提前回收的那段内存会在下一帧被再次创建时拿到错位的地址。本队\ours{修正了缺页处理与引用计数这两处}，连续推理随之稳定，首次推理的检测结果与 Linux 一致。这部分对 GEM 缓冲对象的修复以合并请求 \href{https://github.com/rcore-os/tgoskits/pull/1364}{\#1364} 提交至上游，已合并。

\paragraph{零拷贝底座}
三个加速器各有独立的总线主控通道，若每一级都把上一级的结果先搬回 CPU、再拷贝给下一级，单帧仅在搬运上就要付出与一次推理相当的内存带宽与缓存抖动。让它们读写同一段物理内存是整条流水线能够流动起来的前提。Starry 此前没有对应机制，现已\ours{新增 dma-heap 分配器与 dma-buf 零拷贝导入}：一段物理连续内存经 \texttt{/dev/dma\_heap} 只分配一次，以文件描述符在 JPU、RGA、NPU 之间传递，每个部件凭描述符解析出同一个物理地址，整条 \texttt{MJPEG\,→\,JPU\,→\,RGA\,→\,NPU} 通路因此全程零拷贝。应用以一次 \texttt{DMA\_HEAP\_IOCTL\_ALLOC} 提交一份 \texttt{dma\_heap\_allocation\_data} 请求长度，该结构在 LP64 下恰为 24 字节，ioctl 编号由内核侧按 \texttt{\_IOWR(\textquotesingle H\textquotesingle, 0, ...)} 推算并以编译期断言固定为 \texttt{0xC0184800}，与主线 Linux 的取值逐位一致，任何字段或编号的偏移都会在编译时暴露。分配落在 4\,GiB 以下的物理空间、按页对齐且物理连续，这两条约束各有来由：物理连续是因为 RGA2 与 NPU 在本阶段以设备侧地址翻译部件（IOMMU）旁路、物理直通的方式沿一个起始物理地址顺序访问整段缓冲，中途若被分页打断便会读到错误的页；32 位地址上限则对应 RGA2 把一个裸 32 位物理基址直接写入自身寄存器的约束，分配器因此以 \texttt{u32::MAX} 作为 DMA 掩码并把请求长度限制在 32 位以内。设备物理直通带来缓存一致性问题：\ours{在把缓冲交给引擎之前沿写入设备的方向同步一次、将脏的 CPU 缓存行刷回内存，待引擎完成、读取结果之前再沿读回 CPU 的方向同步一次、使目标缓冲对应的缓存行失效}，这两次方向相反的同步只针对实际参与本次操作的源与目标缓冲，CPU 与设备据此看到一致的内容。

\subsection{在真机上的结果}

在 OrangePi-5-Plus 上按分级测量整条流水线，端到端时延阶梯见推理链路优化一节。

分级来看，RGA 承担的色彩空间转换\keyres{较 CPU 快约 \nRgaSpeedup\,倍}，缩放补边亦接近归零。JPU 单帧解码中位约 1\,ms，解出的帧与 Linux 逐字节相同。推理本身（\texttt{run}）因加速器时钟固定在 1.0\,GHz 不做动态降频，与 Linux 在同一时钟下处于同一水平，取球用的 480×640 模型约 4.7\,ms。
